\section{Tổng kết và Hướng phát triển}

\subsection{Kết quả đạt được}
Nhóm 09 đã hoàn thành bản MVP của hệ thống \textbf{E-Hub} đáp ứng đầy đủ 17 Use Case thiết kế từ Submission 1--3:
\begin{itemize}
    \item Hoàn thiện đầy đủ giao diện Sinh viên (Mobile + Responsive Web) và Dashboard Operator ở cả chế độ Sáng/Tối và song ngữ Việt/Anh (323 khóa).
    \item Hiện thực đầy đủ tầng nghiệp vụ \texttt{services/} và dữ liệu giả lập quy mô 5 Hub, 700 xe, 830 chỗ đỗ, 220 cổng sạc.
    \item Vượt qua toàn bộ \textbf{87/87 test case} tự động trên \texttt{pytest} và chứng minh thành công cả 3 demo bắt buộc (race condition đặt xe, lọc gói trễ, cô lập dữ liệu mô phỏng).
\end{itemize}

\subsection{Bài học kinh nghiệm}
\begin{itemize}
    \item Việc thống nhất hợp đồng dữ liệu (\texttt{DATA\_CONTRACT.md}) ngay từ đầu giúp các thành viên làm giao diện và logic song song mà không bị xung đột tên trường khi tích hợp.
    \item Tách biệt tầng Presentation và Service giúp việc viết kiểm thử tự động (\texttt{pytest}) nhanh chóng và cô lập lỗi hiệu quả.
    \item \textit{\_TODO: Bổ sung thêm bài học kinh nghiệm sau buổi họp tổng kết và quay demo của nhóm.\_}
\end{itemize}

\subsection{Hướng phát triển}
\begin{itemize}
    \item Chuyển đổi dữ liệu giả lập \texttt{seed.json} sang hệ quản trị cơ sở dữ liệu thực tế (PostgreSQL/Redis) và sử dụng WebSocket thay cho cơ chế polling mỗi 5 giây.
    \item Tích hợp xác thực SSO thực tế của ĐHQG-HCM và kết nối luồng telemetry từ thiết bị GPS/IoT gắn trên phương tiện.
\end{itemize}
